\chapter{Proces Poissona: losowe zliczanie zdarzeń}
\label{ch:poisson-process}
\paragraph{Odsyłacze międzytomowe.} Estymację intensywności procesu Poissona i modele Poisson GLM omawia tom IV, rozdziały~\ref{IV-ch:conjugate-families} i~\ref{IV-ch:glm}.


\section{Motywacja: od losowych odstępów do procesu zliczającego}
W poprzednich rozdziałach skonstruowaliśmy procesy jako rodziny
zmiennych losowych, a następnie opisaliśmy informację dostępną w czasie.
Teraz zajmiemy się pierwszym ważnym modelem, w którym \emph{cała
trajektoria} ma prostą interpretację. Liczymy zgłoszenia do serwisu,
przybycia klientów, impulsy z detektora promieniowania albo błędy
transmisji. Interesują nas nie tylko liczby zdarzeń, lecz także czasy
ich wystąpienia, odstępy, zdarzenia po chwili obserwacji i reakcja na
zmianę intensywności.

Dwa pozornie różne opisy mogą prowadzić do tego samego procesu:
\begin{itemize}
\item odstępy między kolejnymi zgłoszeniami są niezależne i wykładnicze;
\item liczby zgłoszeń w rozłącznych przedziałach są niezależne, a liczba
w przedziale długości $t$ ma rozkład Poissona o parametrze $\lambda t$.
\end{itemize}
Udowodnimy równoważność tych opisów, zamiast przyjmować ją jako
nieuzasadnioną własność modelu.

\begin{definicja}[Czasy nadejścia i proces zliczający]
Niech $(\Omega,\calF,\Prob)$ będzie przestrzenią probabilistyczną,
a $E_1,E_2,\ldots$ dodatnimi zmiennymi losowymi. Niech
$T_0:=0$ i $T_n:=\sum_{k=1}^n E_k$, $n\ge1$.
Jeżeli $T_n\to\infty$ prawie na pewno, to dla $t\ge0$ określamy
\[
N_t:=\sum_{n=1}^\infty\ind{\{T_n\le t\}}
      =\max\{n\in\N_0:T_n\le t\}.
\]
Proces $N=(N_t)_{t\ge0}$ jest \emph{procesem zliczającym} zdarzenia
w chwilach $T_1,T_2,\ldots$. Warunek $T_n\to\infty$ oznacza, że na
żadnym ograniczonym przedziale nie występuje nieskończenie wiele
zdarzeń (brak eksplozji).
\end{definicja}

Każda trajektoria $N$ jest niemalejącą funkcją schodkową, prawostronnie
ciągłą i posiadającą granice lewostronne; jej skoki są równe jeden,
gdyż $E_k>0$. Dla przejrzystości wszystkie twierdzenia tego rozdziału
dotyczą czasu $[0,\infty)$.

\begin{center}
\begin{tikzpicture}[>=Latex,x=1.05cm,y=.62cm]
\draw[->] (-.3,0)--(8.1,0) node[right] {$t$};
\draw[->] (0,-.3)--(0,4.8) node[above] {$N_t$};
\foreach \j in {1,2,3,4} {\draw (-.09,\j)--(.09,\j) node[left=4pt] {$\j$};}
\draw[very thick,blue!70!black] (0,0)--(1.05,0)--(1.05,1)--(2.55,1)--(2.55,2)--(3.6,2)--(3.6,3)--(6.2,3)--(6.2,4)--(7.75,4);
\foreach \x/\j in {1.05/1,2.55/2,3.6/3,6.2/4} {
\fill[blue!70!black] (\x,\j) circle (2pt);
\draw[fill=white,blue!70!black] (\x,{\j-1}) circle (2pt);
\draw[densely dashed,gray] (\x,0)--(\x,\j);
\node[below=4pt] at (\x,0) {$T_{\j}$};}
\end{tikzpicture}
\par\smallskip\small Przykładowa trajektoria: $N_t$ liczy zdarzenia w $(0,t]$;
punkty pełne pokazują przyjętą prawostronną ciągłość.
\end{center}

\section{Konstrukcja ze zmiennych wykładniczych}
\begin{definicja}[Jednorodny proces Poissona]
Niech $\lambda>0$. \emph{Jednorodnym procesem Poissona o intensywności
$\lambda$} nazywamy proces zliczający skonstruowany z niezależnych
zmiennych $E_1,E_2,\ldots$ o wspólnym rozkładzie wykładniczym
$\operatorname{Exp}(\lambda)$, czyli
$\Prob(E_k>u)=e^{-\lambda u}$ dla $u\ge0$.
Czasy $T_n$ nazywamy chwilami nadejścia, a $E_n$ odstępami między
nadejściami.
\end{definicja}

Najpierw musimy sprawdzić, że konstrukcja ma sens dla wszystkich
$t\ge0$. Nie wolno poprzestać na intuicji, że ``czasy są dodatnie'':
suma nieskończenie wielu dodatnich liczb może być skończona.

\begin{lemat}[Brak eksplozji dla stałej intensywności]
Jeżeli $(E_n)_{n\ge1}$ są niezależnymi zmiennymi o rozkładzie
$\operatorname{Exp}(\lambda)$, $\lambda>0$, to
$T_n=\sum_{k=1}^n E_k\to\infty$ prawie na pewno.
\end{lemat}
\begin{proof}
Ustalmy $a=1/\lambda$. Wydarzenia $A_n=\{E_n\ge a\}$ są niezależne
i $\Prob(A_n)=e^{-1}>0$. Mamy więc
$\sum_n\Prob(A_n)=\infty$. Drugi lemat Borela--Cantellego daje
$\Prob(A_n\ \text{nieskończenie często})=1$. Na każdym takim
wyniku $\omega$ nieskończenie wiele składników $E_n(\omega)$ jest
nie mniejszych od $a$. Suma częściowa $T_n(\omega)$ jest niemalejąca
i przekracza dowolną ustaloną wielokrotność $a$, a zatem dąży do
nieskończoności.
\end{proof}

\begin{lemat}[Gęstość chwil nadejścia]
Dla $\lambda>0$, $n\ge1$ i niezależnych
$E_1,\ldots,E_n\sim\operatorname{Exp}(\lambda)$ wektor
$(T_1,\ldots,T_n)$ ma gęstość
\[
(s_1,\ldots,s_n)\longmapsto
\lambda^n e^{-\lambda s_n}\ind{\{0<s_1<\cdots<s_n\}}.
\]
W szczególności $T_n$ ma gęstość gamma
$f_{T_n}(t)=\lambda^n t^{n-1}e^{-\lambda t}/(n-1)!$ dla $t>0$.
\end{lemat}
\begin{proof}
Wektor $(E_1,\ldots,E_n)$ ma na dodatnim ortancie gęstość
$\lambda^n\exp(-\lambda\sum_{k=1}^n e_k)$.
Podstawiamy $s_k=e_1+\cdots+e_k$. Odwrotna transformacja ma postać
$e_1=s_1$, $e_k=s_k-s_{k-1}$ dla $k\ge2$ i trójkątny jakobian
równy $1$. Dodatniość wszystkich $e_k$ jest równoważna
$0<s_1<\cdots<s_n$, natomiast suma $e_k$ równa się $s_n$.
Otrzymujemy wskazaną gęstość wspólną. Dla $s_n=t$ całkujemy po
sympleksie $0<s_1<\cdots<s_{n-1}<t$ o objętości
$t^{n-1}/(n-1)!$, co daje gęstość $T_n$.
\end{proof}

\begin{twierdzenie}[Rozkład liczby nadejść i niezależność przyrostów]
\label{thm:poisson-fdd}
Niech $N$ będzie procesem Poissona o intensywności $\lambda>0$.
Wówczas $N_0=0$ oraz dla każdego $t\ge0$ i $n\in\N_0$
\[
\Prob(N_t=n)=e^{-\lambda t}\frac{(\lambda t)^n}{n!}.
\]
Jeżeli $0=t_0<t_1<\cdots<t_m$, to przyrosty
$N_{t_k}-N_{t_{k-1}}$, $k=1,\ldots,m$, są niezależne, a ich
rozkłady to odpowiednio
$\operatorname{Poisson}(\lambda(t_k-t_{k-1}))$.
W szczególności przyrosty są stacjonarne.
\end{twierdzenie}
\begin{proof}
Zdarzenie $\{N_t=n\}$ oznacza $\{T_n\le t<T_{n+1}\}$,
przy $n=0$ rozumiane jako $\{E_1>t\}$.
Dla $n\ge1$ po ustaleniu $0<s_1<\cdots<s_n<t$ prawdopodobieństwo,
że następny niezależny odstęp przekracza $t-s_n$, wynosi
$e^{-\lambda(t-s_n)}$. Wykorzystując poprzedni lemat, otrzymujemy
gęstość wspólną $n$ kolejnych chwil nadejścia \emph{i braku następnego
zdarzenia do chwili $t$}:
\[
\lambda^n e^{-\lambda s_n}e^{-\lambda(t-s_n)}
\ind{\{0<s_1<\cdots<s_n<t\}}
=\lambda^n e^{-\lambda t}\ind{\{0<s_1<\cdots<s_n<t\}}.
\]
Jej całka po sympleksie ma wartość
$e^{-\lambda t}\lambda^n t^n/n!$. Dla $n=0$ wzór wynika
bezpośrednio z rozkładu $E_1$.

Aby udowodnić \emph{niezależność}, nie wystarczy znać każdego
rozkładu brzegowego. Ustalmy liczby całkowite $n_1,\ldots,n_m\ge0$
i $n=\sum_k n_k$. Warunkowo względem $N_{t_m}=n$ chwile
$T_1,\ldots,T_n$ mają na uporządkowanym sympleksie stałą gęstość
$n!/t_m^n$: otrzymujemy ją przez podzielenie wyliczonej przed chwilą
gęstości wspólnej przez $\Prob(N_{t_m}=n)$.
Jest to prawo statystyk pozycyjnych $n$ niezależnych punktów
jednostajnych na $(0,t_m)$. Rzeczywiście, $n$ punktów
nieuporządkowanych ma stałą gęstość $t_m^{-n}$ na
$(0,t_m)^n$, a niemal każdy ich układ ma dokładnie $n!$
różnych uporządkowań. Po uporządkowaniu otrzymujemy więc
gęstość $n!/t_m^n$ na sympleksie. Z kolei przypisanie
$n$ niezależnych punktów do $m$ przedziałów ma
$\frac{n!}{n_1!\cdots n_m!}$ możliwych ustawień
z określonymi licznościami $n_1,\ldots,n_m$; prawdopodobieństwo
każdego ustawienia to iloczyn potęg ich względnych długości.
Dlatego liczby punktów w kolejnych
przedziałach $(t_{k-1},t_k]$ mają warunkowo rozkład wielomianowy:
\[
\Prob\bigl(N_{t_k}-N_{t_{k-1}}=n_k,\ 1\le k\le m
 \mid N_{t_m}=n\bigr)
=\frac{n!}{\prod_k n_k!}\prod_{k=1}^m
 \left(\frac{t_k-t_{k-1}}{t_m}\right)^{n_k}.
\]
Mnożąc tę liczbę przez $e^{-\lambda t_m}(\lambda t_m)^n/n!$
i rozkładając $e^{-\lambda t_m}=\prod_k e^{-\lambda(t_k-t_{k-1})}$,
uzyskujemy
\[
\Prob\bigl(N_{t_k}-N_{t_{k-1}}=n_k,\ 1\le k\le m\bigr)
=\prod_{k=1}^m e^{-\lambda(t_k-t_{k-1})}
  \frac{[\lambda(t_k-t_{k-1})]^{n_k}}{n_k!}.
\]
To dokładnie iloczyn rozkładów brzegowych. Dowód zakończony.
\end{proof}


\begin{twierdzenie}[Dwa równoważne opisy jednorodnego procesu Poissona]
\label{thm:poisson-reverse}
Niech $\lambda>0$, a $N=(N_t)_{t\ge0}$ będzie procesem
o niemalejących trajektoriach c\`adl\`ag, wartościach
w $\N_0$, $N_0=0$ i skokach o wielkości dokładnie $1$.
Załóżmy, że $N$ ma niezależne, stacjonarne przyrosty
oraz $N_t\sim\operatorname{Poisson}(\lambda t)$ dla każdego
$t\ge0$. Wtedy chwile nadejść
$T_k:=\inf\{t\ge0:N_t\ge k\}$, $k\ge1$,
są prawie na pewno skończone, a odstępy
$E_k:=T_k-T_{k-1}$, $T_0=0$, są niezależne i mają
wspólny rozkład $\operatorname{Exp}(\lambda)$.
\end{twierdzenie}
\begin{proof}
\emph{Krok 1: chwile $T_k$ istnieją i są skończone.}
Dla każdego $k\ge1$ oraz $t\ge0$, dzięki niemalejącym
trajektoriom i ich prawostronnej ciągłości,
$\{T_k\le t\}=\{N_t\ge k\}$.
W szczególności $\Prob(T_k>t)=\Prob(N_t<k)$.
Dla ustalonego $k$ prawa strona dąży do zera, gdy
$t\to\infty$: dla $N_t\sim\operatorname{Poisson}(\lambda t)$
na przykład nierówność Czebyszewa daje
$\Prob(N_t<k)\le\lambda t/(\lambda t-k)^2$
dla $\lambda t>k$.
Zatem każdy $T_k$ jest skończony prawie na pewno;
przecinamy te przeliczalnie liczne zdarzenia.

\emph{Krok 2: rozkłady wszystkich skończonych rodzin czasów.}
Prawo wektora $(T_1,\ldots,T_m)$ jest wyznaczone
przez prawdopodobieństwa
\[
\Prob(T_1\le t_1,\ldots,T_m\le t_m)
=\Prob(N_{t_1}\ge1,\ldots,N_{t_m}\ge m),
\qquad t_1,\ldots,t_m\ge0.
\]
Prawa strona zależy wyłącznie od rozkładów
skończeniewymiarowych $N$. Te zaś są jednoznacznie
wyznaczone przez $N_0=0$ oraz rozkłady niezależnych
przyrostów: dla uporządkowanych chwil $0=t_0<t_1<\cdots<t_r$
wektor wartości $N_{t_j}$ jest obrazem wektora
niezależnych przyrostów $N_{t_j}-N_{t_{j-1}}$ przez
odwzorowanie polegające na sumowaniu współrzędnych.

\emph{Krok 3: porównanie z wcześniej skonstruowanym modelem.}
W twierdzeniu~\ref{thm:poisson-fdd} skonstruowaliśmy
proces $\widehat N$ z niezależnych odstępów
$\widehat E_k\sim\operatorname{Exp}(\lambda)$ i wykazaliśmy,
że ma dokładnie takie same rozkłady skończeniewymiarowe jak $N$.
Krok 2 oznacza zatem, że dla każdego $m$
$(T_1,\ldots,T_m)$ i $(\widehat T_1,\ldots,\widehat T_m)$
mają te same rozkłady. Przekształcenie liniowe
$(t_1,\ldots,t_m)\mapsto(t_1,t_2-t_1,\ldots,t_m-t_{m-1})$
pokazuje identyczność praw wektorów
$(E_1,\ldots,E_m)$ i $(\widehat E_1,\ldots,\widehat E_m)$.
Drugi wektor ma rozkład produktowy
$\operatorname{Exp}(\lambda)^{\otimes m}$.
Ponieważ $m$ jest dowolne, wszystkie odstępy $E_k$
są niezależne i jednakowo wykładnicze, co kończy dowód.
\end{proof}

\begin{uwaga}
Kierunek odwrotny nie wymagał dodatkowego założenia o
mocnej własności Markowa w losowych chwilach $T_k$.
Najpierw porównaliśmy \emph{rozkłady skończeniewymiarowe},
a dopiero potem wyprowadziliśmy z nich wspólne prawo
czasów nadejść. To dobry przykład konstrukcyjnego
użycia twierdzenia o jednoznaczności prawa procesu.
\end{uwaga}

Wynika stąd $\E N_t=\Var(N_t)=\lambda t$,
a dla $0\le s\le t$ mamy
$\Cov(N_s,N_t)=\Var(N_s)=\lambda s$: zapisujemy
$N_t=N_s+(N_t-N_s)$ i korzystamy z niezależności przyrostu.

\begin{uwaga}[Co dokładnie znaczy niezależność przyrostów?]
Własność oznacza niezależność \emph{liczb zdarzeń} na rozłącznych
przedziałach, nie niezależność $N_s$ i $N_t$. Przeciwnie,
$N_t\ge N_s$ dla $t\ge s$ i zwykle
$\Cov(N_s,N_t)=\lambda s>0$. Co więcej, dla ustalonych
$t\ge0$ i $h>0$ mamy
$\Prob(N_{t+h}-N_t=0)=e^{-\lambda h}$,
$\Prob(N_{t+h}-N_t=1)=\lambda h+O(h^2)$
oraz $\Prob(N_{t+h}-N_t\ge2)=O(h^2)$ przy $h\downarrow0$.
To wyjaśnia pojęcie \emph{intensywności}: prawdopodobieństwo
jednego nadejścia na bardzo krótkim przedziale jest w pierwszym
przybliżeniu równe $\lambda h$.
\end{uwaga}

\section{Własność braku pamięci, filtracja i martyngał}
\begin{stwierdzenie}[Brak pamięci rozkładu wykładniczego]
Jeżeli $E\sim\operatorname{Exp}(\lambda)$, $\lambda>0$, to
$\Prob(E>s+t\mid E>s)=\Prob(E>t)$ dla wszystkich $s,t\ge0$.
Wśród rozkładów na $[0,\infty)$ o ciągłej funkcji przeżycia,
ściśle dodatniej dla argumentów skończonych i przyjmującej wartość
$1$ w zerze, ta własność charakteryzuje rozkład wykładniczy.
\end{stwierdzenie}
\begin{proof}
Pierwsza równość wynika ze wzoru
$e^{-\lambda(s+t)}/e^{-\lambda s}=e^{-\lambda t}$.
Dla odwrotnego kierunku oznaczmy $g(t)=\Prob(E>t)$.
Z własności warunkowania wynika $g(s+t)=g(s)g(t)$.
Funkcja $h(t):=-\log g(t)$ jest ciągła, spełnia równanie
$h(s+t)=h(s)+h(t)$ dla $s,t\ge0$ i $h(0)=0$.
Stąd najpierw dla argumentów wymiernych, a potem przez ciągłość
dla wszystkich $t\ge0$ otrzymujemy $h(t)=\lambda t$.
Ponieważ $g$ jest nierosnąca, $\lambda\ge0$; aby otrzymać
właściwy rozkład ze skończoną wartością $E$ prawie na pewno,
musimy mieć $\lim_{t\to\infty}g(t)=0$, a więc $\lambda>0$.
\end{proof}

\begin{twierdzenie}[Proces Poissona jest procesem Markowa]
Niech $N$ będzie procesem Poissona o intensywności $\lambda>0$,
a $\calF_t^N:=\sigma(N_s:0\le s\le t)$ jego naturalną
filtracją. Dla $0\le s<t$ i każdego $k\in\N_0$ zachodzi
\[
\Prob(N_t=N_s+k\mid\calF_s^N)
=e^{-\lambda(t-s)}\frac{[\lambda(t-s)]^k}{k!}
\quad\text{prawie na pewno}.
\]
W szczególności rozkład przyszłego stanu warunkowo względem
całej przeszłości zależy wyłącznie od aktualnego $N_s$.
\end{twierdzenie}
\begin{proof}
Z twierdzenia~\ref{thm:poisson-fdd} wynika niezależność
$N_t-N_s$ od wszystkich wektorów
$(N_{u_1},\ldots,N_{u_r})$ z $u_j\le s$:
możemy każdy $N_{u_j}$ zapisać jako sumę przyrostów na podziale
przedziału $[0,s]$ zawierającym wszystkie $u_j$.
Klasa zdarzeń cylindrycznych wyznaczonych przez takie wektory
jest $\pi$-układem generującym $\calF_s^N$.
Przez twierdzenie o jednoznaczności miar (argument
$\pi$--$\lambda$) niezależność rozciąga się na całe
$\calF_s^N$. Zatem rozkład warunkowy przyrostu jest jego
rozkładem bezwarunkowym, wyliczonym wcześniej.
Ponieważ $N_t=N_s+(N_t-N_s)$, dostajemy tezę.
\end{proof}

\begin{wniosek}[Kompensowany proces Poissona]
Dla procesu Poissona o intensywności $\lambda>0$ proces
$M_t:=N_t-\lambda t$ jest kwadratowalnym martyngałem względem
filtracji naturalnej $\mathbb F^N$, a
$\E(M_t-M_s)^2=\lambda(t-s)$ dla $0\le s\le t$.
\end{wniosek}
\begin{proof}
Dla $s\le t$ niezależność przyrostów daje
$\E[N_t-N_s\mid\calF_s^N]=\lambda(t-s)$,
więc $\E[M_t\mid\calF_s^N]=M_s$.
Całkowalność i kwadratowalność wynikają z momentów
rozkładu Poissona. Wreszcie
$\E(M_t-M_s)^2=\Var(N_t-N_s)=\lambda(t-s)$.
\end{proof}


\begin{center}
\begin{tikzpicture}[>=Latex,x=.95cm,y=.78cm]
\draw[->] (0,0)--(10,0) node[right] {$t$};
\foreach \x/\lab in {1/$T_1$,3/$T_2$,6/$T_3$,8/$T_4$}{
  \draw[red!65!black,very thick] (\x,-.16)--(\x,.16);
  \node[above=3pt] at (\x,.17) {\lab};
}
\draw[<->,blue!60!black] (1,-.52)--(3,-.52)
 node[midway,below] {$E_2=T_2-T_1$};
\draw[<->,blue!60!black] (3,-1.1)--(6,-1.1)
 node[midway,below] {$E_3=T_3-T_2$};
\node[align=left] at (5,1.05)
 {czasy nadejść: $T_1,T_2,\ldots$\qquad odstępy: $E_1,E_2,\ldots$};
\end{tikzpicture}
\par\smallskip\small Oś chwil nadejścia. Zliczenia na przedziale $(s,t]$
 zależą od liczby zaznaczonych punktów, a odstępy opisują
 czas pomiędzy kolejnymi punktami.
\end{center}

\section{Czasy nadejść i warunkowanie na liczbie zdarzeń}
\begin{twierdzenie}[Chwila $n$-tego zdarzenia]
Dla $n\ge1$ w procesie Poissona z intensywnością $\lambda>0$
czas $T_n$ ma rozkład gamma o parametrach $n$ i $\lambda$,
przy konwencji, że drugi parametr jest intensywnością. Zachodzi
\[
\Prob(T_n>t)=e^{-\lambda t}\sum_{k=0}^{n-1}
\frac{(\lambda t)^k}{k!},\qquad
\E T_n=\frac n\lambda,\qquad \Var(T_n)=\frac n{\lambda^2}.
\]
\end{twierdzenie}
\begin{proof}
Gęstość została wyprowadzona wcześniej. Zdarzenie
$\{T_n>t\}$ jest identyczne ze zdarzeniem $\{N_t\le n-1\}$;
stąd wzór na funkcję przeżycia. Ponieważ
$T_n=E_1+\cdots+E_n$, niezależność i momenty rozkładu
wykładniczego dają wskazane średnią i wariancję.
\end{proof}

\begin{twierdzenie}[Jednostajność chwil nadejść przy ustalonej liczbie]
\label{thm:uniform-poisson-times}
Niech $t>0$ i $n\ge1$. Warunkowo względem zdarzenia
$\{N_t=n\}$ wektor $(T_1,\ldots,T_n)$ ma taki sam rozkład,
jak uporządkowane statystyki pozycyjne $n$ niezależnych
zmiennych o rozkładzie jednostajnym na $(0,t)$.
W szczególności dla $0<s<t$ i $0\le k\le n$
\[
\Prob(N_s=k\mid N_t=n)
=\binom nk\left(\frac{s}{t}\right)^k
\left(1-\frac{s}{t}\right)^{n-k}.
\]
\end{twierdzenie}
\begin{proof}
W dowodzie twierdzenia~\ref{thm:poisson-fdd} otrzymaliśmy
warunkową gęstość $n!/t^n$ na sympleksie
$0<s_1<\cdots<s_n<t$. Dokładnie taka jest gęstość
statystyk pozycyjnych z jednostajnego losowania
(gęstość wspólna nieuporządkowanego wektora to $t^{-n}$,
a każdy porządek z $n!$ ma jednakowe prawdopodobieństwo).
Liczba wylosowanych punktów należących do $(0,s]$ jest sumą
$n$ niezależnych zmiennych Bernoulliego z parametrem $s/t$,
co daje wzór dwumianowy.
\end{proof}

Warto zauważyć zmianę perspektywy: bez warunkowania liczby
nadejść w rozłącznych przedziałach są niezależne, ale po warunku
$N_t=n$ przestają być niezależne, bo ich suma musi wynosić $n$.

\begin{zadanie}[Cztery zgłoszenia w godzinę]
Zgłoszenia pojawiają się zgodnie z procesem Poissona o
intensywności $\lambda=4$ na godzinę. Obliczyć: (a)
prawdopodobieństwo dokładnie trzech zgłoszeń w pierwszej godzinie;
(b) prawdopodobieństwo, że do chwili $15$ minut nadejdzie
przynajmniej jedno zgłoszenie; (c) prawdopodobieństwo,
że przy dokładnie czterech zgłoszeniach w godzinie dwa z nich
wystąpiły w pierwszych piętnastu minutach;
(d) oczekiwany czas czwartego zgłoszenia.

\emph{Rozwiązanie.} Przyjmujemy czas w godzinach.
(a) Z rozkładu Poissona dostajemy
$\Prob(N_1=3)=e^{-4}4^3/3!$.
(b) Na odcinku $[0,1/4]$ oczekiwana liczba zgłoszeń wynosi
$1$, więc $\Prob(N_{1/4}\ge1)=1-e^{-1}$.
(c) Z twierdzenia~\ref{thm:uniform-poisson-times} liczba
zgłoszeń w pierwszym kwadransie pod warunkiem $N_1=4$
ma rozkład $\operatorname{Bin}(4,1/4)$; prawdopodobieństwo
wynosi $\binom42(1/4)^2(3/4)^2=27/128$.
(d) Suma czterech niezależnych odstępów wykładniczych ma średnią
$4/\lambda=1$ godzinę. Średnia ta nie oznacza oczywiście, że
czwarte zgłoszenie przychodzi dokładnie po godzinie.
\end{zadanie}


\begin{zadanie}[Dokładne chwile nadejść po warunkowaniu]
Klienci przybywają według procesu Poissona o intensywności
$\lambda>0$. Wiemy, że w ciągu dwóch godzin przybyło dokładnie
dwóch klientów. Niech $T_1<T_2$ będą ich chwilami przybycia.
Obliczyć $\Prob(T_1<1<T_2\mid N_2=2)$,
$\E[T_1\mid N_2=2]$ oraz gęstość wspólną $(T_1,T_2)$
przy tym warunku.

\emph{Rozwiązanie.} Warunkowo względem $N_2=2$
dwie chwile są statystykami pozycyjnymi dwóch niezależnych
zmiennych o rozkładzie jednostajnym na $[0,2]$.
Ich uporządkowana gęstość wynosi
$f(s_1,s_2)=\frac{2!}{2^2}=\frac12$ na trójkącie
$0<s_1<s_2<2$ i jest zerowa poza nim. Zdarzenie
$T_1<1<T_2$ oznacza, że jedna chwila znajduje się
w $(0,1)$, a druga w $(1,2)$; dlatego ma warunkowe
prawdopodobieństwo $2(1/2)^2=1/2$.
Dystrybuanta minimum dwóch zmiennych jednostajnych daje
$\Prob(T_1>u\mid N_2=2)=(1-u/2)^2$ dla $0\le u\le2$.
Stąd, całkując funkcję przeżycia,
\[
\E[T_1\mid N_2=2]
=\int_0^2(1-u/2)^2\,\dd u=\frac23.
\]
Wyniki nie zależą od $\lambda$: po ustaleniu liczby nadejść
intensywność nie wpływa na względne położenie ich chwil.
\end{zadanie}

\section{Superpozycja i rozrzedzanie}
\begin{twierdzenie}[Superpozycja]
Niech $N^{(1)},\ldots,N^{(r)}$ będą niezależnymi procesami
Poissona o intensywnościach $\lambda_1,\ldots,\lambda_r>0$.
Wtedy $N_t:=\sum_{j=1}^rN_t^{(j)}$ jest procesem Poissona
o intensywności $\lambda:=\sum_j\lambda_j$.
\end{twierdzenie}
\begin{proof}
Trajektorie sumy są niemalejące i prawostronnie ciągłe.
W ustalonej chwili prawdopodobieństwo jednoczesnego skoku
dwóch niezależnych procesów wynosi zero: warunkowo względem
skończonej liczby skoków pierwszego procesu na przedziale
$[0,m]$, ich chwile mają rozkład z gęstością;
dla każdej deterministycznej chwili prawdopodobieństwo skoku
drugiego procesu jest zero. Sumując po chwilach i po $m\in\N$,
wykluczamy wspólne skoki prawie na pewno.

Dla rozłącznych przedziałów przyrosty sumy są niezależne,
ponieważ wszystkie przyrosty procesów składowych z tych
przedziałów tworzą niezależną rodzinę. Dla pojedynczego
przedziału długości $h$ funkcja tworząca prawdopodobieństwa
przyrostu wynosi
\[
\E z^{\sum_j (N_{t+h}^{(j)}-N_t^{(j)})}
=\prod_j e^{\lambda_j h(z-1)}
=e^{(\sum_j\lambda_j)h(z-1)},\qquad 0\le z\le1.
\]
Jest to funkcja tworząca rozkładu
$\operatorname{Poisson}(\lambda h)$; jednoznaczność
współczynników szeregu potęgowego daje taki właśnie rozkład.
Zatem suma spełnia charakterystykę z twierdzenia
\ref{thm:poisson-fdd}.
\end{proof}

\begin{twierdzenie}[Niezależne rozrzedzanie]
Niech $N$ będzie procesem Poissona o intensywności $\lambda>0$.
Każdy jego skok, niezależnie od innych skoków i od $N$,
oznaczmy jako typu $A$ z prawdopodobieństwem $p\in[0,1]$,
a jako typu $B$ z prawdopodobieństwem $1-p$.
Niech $N_t^A$ oraz $N_t^B$ zliczają odpowiednie typy do chwili $t$.
Wtedy $N^A$ i $N^B$ są \emph{niezależnymi} procesami Poissona
o intensywnościach $\lambda p$ i $\lambda(1-p)$
(z konwencją, że intensywność $0$ oznacza proces zerowy).
\end{twierdzenie}
\begin{proof}
Na ustalonym przedziale długości $h$, warunkowo względem
$N_{t+h}-N_t=n$, liczba oznaczeń $A$ ma rozkład
$\operatorname{Bin}(n,p)$, a liczba $B$ jest dopełnieniem do $n$.
Dla $|u|,|v|\le1$ wspólna funkcja tworząca wynosi
\[
\E\left[u^{N^A_{t+h}-N^A_t}v^{N^B_{t+h}-N^B_t}\right]
=\E[ (pu+(1-p)v)^{N_{t+h}-N_t}]
=e^{\lambda h p(u-1)}e^{\lambda h(1-p)(v-1)}.
\]
Iloczyn wskazuje jednocześnie rozkłady Poissona i
niezależność dwóch liczników na tym przedziale.
Ponieważ podstawowy proces ma niezależne przyrosty, a
oznaczenia na rozłącznych przedziałach są niezależne,
rozumowanie stosuje się łącznie do dowolnej skończonej
rodziny rozłącznych przedziałów. Tym samym oba powstałe
procesy są niezależne jako rodziny zmiennych losowych,
a nie tylko dla każdego ustalonego $t$.
\end{proof}

\begin{zadanie}[Dwa strumienie i klasyfikacja zgłoszeń]
Przychodzą niezależne zgłoszenia z działów $A$ i $B$
o intensywnościach odpowiednio $2$ i $3$ na godzinę.
Dodatkowo każde zgłoszenie, niezależnie od pozostałych, jest
pilne z prawdopodobieństwem $1/5$. Obliczyć
prawdopodobieństwo braku pilnych zgłoszeń w ciągu dwóch godzin;
oraz prawdopodobieństwo, że pierwsze zgłoszenie całego systemu
pochodzi z działu $A$.

\emph{Rozwiązanie.} Superpozycja daje intensywność $5$.
Rozrzedzenie pilnych zgłoszeń ma intensywność $5/5=1$;
w ciągu dwóch godzin liczba pilnych zgłoszeń jest więc
$\operatorname{Poisson}(2)$, a prawdopodobieństwo zera wynosi
$e^{-2}$.
Pierwsze czasy nadejść z działów $A$ i $B$ są niezależne,
$U\sim\operatorname{Exp}(2)$ i
$V\sim\operatorname{Exp}(3)$. Zatem
\[
\Prob(U<V)=\int_0^\infty \Prob(V>u)\,2e^{-2u}\,\dd u
=\int_0^\infty 2e^{-5u}\,\dd u=\frac25.
\]
Interpretacja: spośród wszystkich zdarzeń frakcja $2/5$
pochodzi średnio z działu $A$, ale wynikiem jest tu
konkretne prawdopodobieństwo typu \emph{pierwszego} zgłoszenia.
\end{zadanie}

\section{Złożony proces Poissona i kompensacja}
\begin{definicja}[Złożony proces Poissona]
Niech $N$ będzie procesem Poissona o intensywności $\lambda>0$,
a $(Y_k)_{k\ge1}$ ciągiem niezależnych rzeczywistych
zmiennych losowych o wspólnym rozkładzie, niezależnym od $N$.
Definiujemy $S_t=\sum_{k=1}^{N_t}Y_k$,
gdzie pusta suma wynosi zero. Proces $S$ nazywamy
\emph{złożonym procesem Poissona}. Skok $S$ w chwili $T_k$
ma wielkość $Y_k$.
\end{definicja}

Model ma szerokie zastosowania: wypadki występują w chwilach
Poissona, natomiast kwota szkody jest niezależnym losowym
\emph{znacznikiem} zdarzenia. Dla $Y_k=1$ odzyskujemy $N_t$.

\begin{twierdzenie}[Moment, wariancja i funkcja charakterystyczna]
Załóżmy w poprzedniej definicji $\E|Y_1|<\infty$.
Wtedy $\E S_t=\lambda t\E Y_1$.
Jeśli ponadto $\E Y_1^2<\infty$, to
$\Var(S_t)=\lambda t\E Y_1^2$.
Bez założeń momentowych zachodzi dla $u\in\R$
\[
\E e^{iuS_t}=\exp\!\bigl(\lambda t(\E e^{iuY_1}-1)\bigr).
\]
Przy $\E|Y_1|<\infty$ proces
$S_t-\lambda t\E Y_1$ jest martyngałem względem
filtracji generowanej przez zaznaczone zdarzenia do chwili $t$.
\end{twierdzenie}
\begin{proof}
Warunkując względem $N_t=n$ i używając niezależności znaczników,
uzyskujemy $\E[S_t\mid N_t=n]=n\E Y_1$,
$\Var(S_t\mid N_t=n)=n\Var(Y_1)$ i
$\E[e^{iuS_t}\mid N_t=n]=(\E e^{iuY_1})^n$.
Po uśrednieniu pierwsza równość jest oczywista. Ze wzoru
na wariancję całkowitą wynika
\begin{align*}
\Var(S_t)&=\E[N_t\Var(Y_1)]
 +\Var(N_t\E Y_1)\\
&=\lambda t\bigl(\Var(Y_1)+(\E Y_1)^2\bigr)
=\lambda t\E Y_1^2.
\end{align*}
Uśrednienie funkcji charakterystycznej wykorzystuje szereg
$\sum_{n\ge0}e^{-\lambda t}(\lambda t)^n z^n/n!
=e^{\lambda t(z-1)}$.
Przyrosty $S_t-S_s$ są niezależne od informacji do chwili $s$
i mają średnią $\lambda(t-s)\E Y_1$:
wynika to z niezależnych przyrostów procesu Poissona i z faktu,
że znaczniki przyszłych zdarzeń są niezależne od przeszłych.
Stąd właściwość martyngałowa po odjęciu deterministycznej średniej.
\end{proof}

\section{Niejednorodny proces Poissona}
Stała intensywność jest wygodna, ale liczba zgłoszeń w ciągu
dnia na ogół zależy od godziny. Następująca konstrukcja pokazuje,
że zmienna intensywność nie wymaga zmiany podstawowych idei.

\begin{definicja}[Proces Poissona z deterministyczną intensywnością]
Niech $a:[0,\infty)\to[0,\infty)$ będzie funkcją borelowską,
całkowalną na każdym przedziale $[0,t]$. Oznaczmy
$A(t)=\int_0^t a(s)\,\dd s$.
Jeżeli $\widetilde N$ jest standardowym procesem Poissona
o intensywności $1$, to proces $N_t=\widetilde N_{A(t)}$
nazywamy \emph{niejednorodnym procesem Poissona o intensywności $a$}.
\end{definicja}

\begin{stwierdzenie}
Dla procesu z poprzedniej definicji i $0\le s<t$ przyrost
$N_t-N_s$ ma rozkład
$\operatorname{Poisson}(\int_s^t a(u)\,\dd u)$, a przyrosty
na rozłącznych przedziałach są niezależne. Ponadto
$\Prob(N_t-N_s=0)=\exp(-\int_s^t a(u)\,\dd u)$.
\end{stwierdzenie}
\begin{proof}
Funkcja $A$ jest ciągła i niemalejąca, a
$A(t)-A(s)=\int_s^t a(u)\,\dd u$.
Stąd $N_t-N_s=\widetilde N_{A(t)}-\widetilde N_{A(s)}$.
Przedziały obrazowe $[A(s),A(t)]$ odpowiadające rozłącznym
przedziałom czasu mają rozłączne wnętrza; przedziały o
zerowej długości dają zerowe przyrosty. Zatem po zastosowaniu
niezależnych i stacjonarnych przyrostów $\widetilde N$
dostajemy wszystkie tezy.
\end{proof}

\begin{przyklad}[Intensywność sezonowa]
Niech $a(t)=2+t$ dla $0\le t\le2$ i $a(t)=4$ dla $t>2$.
W ciągu pierwszej godziny średnia liczba zdarzeń wynosi
$A(1)=\int_0^1(2+s)\,\dd s=5/2$,
a prawdopodobieństwo braku zgłoszeń jest $e^{-5/2}$.
Pomiędzy godziną pierwszą a drugą średnia przyrostu
wynosi $A(2)-A(1)=7/2$. Obie liczby zdarzeń są niezależne,
choć ich rozkłady nie są jednakowe: przyrosty pozostają
niezależne, ale przestają być stacjonarne.
\end{przyklad}

\section{Charakterystyka przez przyrosty i zdarzenia rzadkie}
Dotychczas budowaliśmy proces Poissona z czasów wykładniczych.
Można jednak rozpocząć od modelowania zliczeń w krótkich
przedziałach. Poniższy wynik uzasadnia tę drogę konstrukcji.

\begin{twierdzenie}[Charakterystyka procesu Poissona przez małe przyrosty]
\label{thm:poisson-infinitesimal}
Niech $N=(N_t)_{t\ge0}$ będzie procesem o wartościach
w $\N_0$, $N_0=0$, z niezależnymi i stacjonarnymi
nieujemnymi przyrostami. Załóżmy, że istnieje $\lambda>0$
takie, iż przy $h\downarrow0$
\[
\Prob(N_h=1)=\lambda h+o(h),\qquad
\Prob(N_h\ge2)=o(h).
\]
Wtedy dla każdego $t\ge0$ zachodzi
$N_t\sim\operatorname{Poisson}(\lambda t)$,
a rozkłady skończeniewymiarowe $N$ są takie same jak
rozkłady procesu Poissona o intensywności $\lambda$.
Jeżeli ponadto $N$ ma niemalejące trajektorie c\`adl\`ag,
to proces ten ma prawo procesu Poissona także jako
zmienna losowa w przestrzeni trajektorii c\`adl\`ag.
\end{twierdzenie}
\begin{proof}
Ustalmy $z\in[0,1]$ i funkcję tworzącą
$G_t(z)=\E z^{N_t}$, przyjmując $0^0=1$.
Z niezależnych i stacjonarnych przyrostów wynika
$G_{t+h}(z)=G_t(z)G_h(z)$. Ponieważ
$\Prob(N_h=0)=1-\lambda h+o(h)$, a wkład stanów
$N_h\ge2$ do $G_h(z)$ ma moduł nie większy niż
$\Prob(N_h\ge2)=o(h)$, otrzymujemy
\[
G_h(z)=1+\lambda h(z-1)+o(h).
\]
Zatem dla $t\ge0$ pochodna prawostronna spełnia
\[
\frac{\partial}{\partial t^+}G_t(z)
=\lim_{h\downarrow0}G_t(z)\frac{G_h(z)-1}{h}
=\lambda(z-1)G_t(z).
\]
Dodatkowo z $\Prob(N_h=0)\to1$ oraz niezależnych przyrostów
wynika ciągłość $t\mapsto G_t(z)$: z prawej na mocy
$G_{t+h}=G_tG_h$, a z lewej, dla $0<h\le t$,
$G_t=G_{t-h}G_h$ i $G_h\to1$; ponieważ
$0\le G_{t-h}\le1$, różnica dąży do zera.
Funkcja ciągła z ciągłą pochodną prawostronną zadaną
powyższym równaniem ma jego zwykłe rozwiązanie,
co można sprawdzić dla
$H_t=e^{-\lambda(z-1)t}G_t(z)$:
jej pochodna prawostronna znika w każdym punkcie,
a więc $H$ jest stała (wynika to z twierdzenia
o wartości średniej w wersji dla pochodnych jednostronnych,
lub z zastosowania argumentu o ekstremum do
$H_t\pm\varepsilon t$). Ponieważ $G_0(z)=1$,
\[
G_t(z)=e^{\lambda t(z-1)}.
\]
Współczynniki szeregu potęgowego są równe
$e^{-\lambda t}(\lambda t)^n/n!$, co daje rozkład Poissona.
Ponieważ przyrosty są niezależne i stacjonarne,
wszystkie rozkłady skończeniewymiarowe są wyznaczone
przez właśnie otrzymane rozkłady pojedynczych przyrostów.
Dla procesów c\`adl\`ag na przestrzeni $D([0,\infty),\N_0)$
rozkłady skończeniewymiarowe wyznaczają prawo procesu
(jego borelowską $\sigma$-algebrę generują współrzędne),
co kończy dowód.
\end{proof}

\begin{przyklad}[Rozrzedzony schemat Bernoulliego]
Ustalmy $t>0$ i rozważmy $m$ niezależnych prób
na siatce $t/m,2t/m,\ldots,t$, z prawdopodobieństwem
sukcesu $\lambda t/m$ w każdej próbie (dla dostatecznie
dużych $m$). Liczba sukcesów ma rozkład
$\operatorname{Bin}(m,\lambda t/m)$.
Dla ustalonego $k\ge0$ jej prawdopodobieństwo wynosi
\[
\binom mk\left(\frac{\lambda t}{m}\right)^k
\left(1-\frac{\lambda t}{m}\right)^{m-k}
\longrightarrow e^{-\lambda t}\frac{(\lambda t)^k}{k!}.
\]
Granica wynika z
$\binom mk/m^k\to1/k!$ i
$(1-\lambda t/m)^m\to e^{-\lambda t}$.
Dla wszystkich rozłącznych odcinków siatki liczby
sukcesów są niezależne; analogicznie otrzymujemy
zbieżność ich wspólnych rozkładów do niezależnych
przyrostów Poissona. Jest to matematyczne wyjaśnienie,
dlaczego rozkład Poissona naturalnie opisuje wiele
niezależnych, rzadkich zdarzeń.
\end{przyklad}

\section{Co wynika z modelu, a czego model nie mówi?}
Rozkład Poissona nie jest uniwersalnym opisem każdego
zliczania. Jeśli nadejście klienta zwiększa lub zmniejsza
szansę kolejnego nadejścia, niezależność przyrostów jest
nieadekwatna. Jeśli urządzenie ma okres refrakcji, to
bardzo krótkie odstępy między zdarzeniami nie są wykładnicze.
Na dalszych stronach uogólnimy ideę zliczania przez
\emph{intensywności zależne od stanu}. Otrzymamy łańcuchy
Markowa w czasie ciągłym, których czasy oczekiwania będą
nadal wykładnicze, ale ich parametry będą zmieniać się
po każdym skoku procesu.
% ================================================================
% TOM II. ROZDZIAŁ 5
% ================================================================
